% SOURCE_AUTHORITY: sga5_fr_workpass.tex, lines 9248--9290 (LF-normalized unit).
% SOURCE_SHA256: 791F4EFFC5E02832D5D77ED03518C8156D6F07E4C8238B03545DB93D883FBB28
\subsubsection*{2.2}
Sejam $S$ um esquema e $E$ uma $\mathcal O_S$-módulo localmente livre de posto $r+1$. Denota-se por
\[
P=P_S(E)\xrightarrow{p}S
\]
O fibrado projetivo correspondente. O canônico $\mathcal O_P$-módulo $\mathcal O_P(1)$ define, como é conhecido, graças à sequência exata de Kummer, um elemento
\[
\xi\in H^2(P,\mu)
\]
e denotar-se-á por $\eta$ a imagem de $\xi$ pelo morfismo canônico
\[
H^2(P,\mu)\longrightarrow H^0(S,\R^2p_*(\mu))
\]
proveniente da sucessão espectral de Leray
\[
H^i(S,\R^jp_*(\mu))\Rightarrow H^{i+j}(P,\mu).
\]
A classe $\xi$ identifica-se com uma flecha $A_P\to\mu_P[2]$ de $D(P)$, que define por transposição $\mu_P^{\otimes -1}[-2]\to A_P$, de onde, pela adjunção habitual, obtém-se uma flecha
\[
c:\mu_S^{\otimes -1}[-2]\longrightarrow \R p_*(A_P),
\]
que corresponda a $\eta$ na cohomologia de grau $2$.

Por produto tensorial, a flecha $c$ define, graças a (2.2.1), algumas flechas
\[
c_i:\mu_S^{\otimes -i}[-2i]\longrightarrow \R p_*(A_P),
\]
para todo inteiro $i\ge1$, e convém denotar por $c_0$ a flecha de adjunção
\[
A_S\longrightarrow \R p_*(A_P).
\]
Define-se Assim uma flecha
\[
\gamma=\bigoplus_{i=0}^{r}c_i:
\bigoplus_{i=0}^{r}\mu_S^{\otimes -i}[-2i]\longrightarrow \R p_*(A_P)
\]
de $D(S)$.

\begin{statement}{Teorema 2.2.1}
A frecha $\gamma$ é um isomorfismo.
\end{statement}

Uma vez definida a seta $\gamma$, (2.2.1) é igual à seguinte corolário.
